\chapter{Spectral Reformulation and the Real-Zero Mechanism}
\label{chap:spectral}

\section{Hilbert--P\'olya logic}

In the centred coordinate $s=1/2+iz$, RH is equivalent to all zeros of $\XiR(z)$ being real. The strongest familiar mechanism for real spectral values is self-adjointness.

\begin{theorem}[Conditional spectral theorem]
\label{thm:spectral-conditional}
Suppose there exists a self-adjoint operator $H$ on a Hilbert space and an entire spectral determinant $D(z)$ such that
\begin{equation}
 \XiR(z)=e^{a+bz}D(z),
 \label{eq:spectral-factorization}
\end{equation}
where the exponential factor is non-zero and the zeros of $D$ are exactly the eigenvalues of $H$, with multiplicity. Then RH holds.
\end{theorem}
\begin{proof}
The spectrum of a self-adjoint operator is real. The non-vanishing exponential does not alter zeros, so every zero of $\XiR$ is real. By \cref{eq:rh-real-xi-zeros}, every non-trivial zeta zero lies on $\Re s=1/2$.
\end{proof}

The theorem captures the Hilbert--P\'olya implication. Many quantum-chaotic analogies and operator proposals are motivated by zero statistics \citep{montgomery1973,odlyzko1987,berrykeating1999,keatingsnaith2000}. No operator satisfying \cref{eq:spectral-factorization} is constructed in this chapter.

\section{How the half enters spectrally}

Self-adjointness does not derive the number $1/2$ from scratch. The functional equation first determines the centred coordinate. Specifically,
\[
 z=-i(s-1/2)
\]
converts the conjugate-affine involution $\J(s)=1-\bar s$ into ordinary conjugation. In the intermediate centered coordinate $u=s-1/2$, the action is the genuinely anti-linear map $u\mapsto-\bar u$. Thus the half is the affine shift that centers the symmetry before a spectral reality condition can be imposed.

The exact logical separation is
\begin{align*}
 \text{functional symmetry}&\longrightarrow\text{centre }1/2,\\
 \text{self-adjointness, if constructed}&\longrightarrow\text{reality of the spectral variable}.
\end{align*}

\section{Two-level compression of a spectral problem}

Suppose a candidate operator $H$ carries an involution or grading that decomposes the Hilbert space into two complementary sectors. Near an eigenvalue, an effective two-dimensional subspace may be spanned by sector components. Under additional hypotheses, a determinant condition can reduce to cancellation of two amplitudes.

\textbf{INTERPRETACJA.} In such a construction the information-spinor model could be read as an effective two-level description of a richer spectral object rather than as a claim that the entire zeta function is literally a qubit amplitude. This interpretation is not used as a theorem premise.

\section{A $2\times2$ toy Hamiltonian}

Consider the Hermitian matrix
\begin{equation}
 H(x,t)=
 \begin{pmatrix}
  x & q(t)\\
  \overline{q(t)} & -x
 \end{pmatrix},
 \label{eq:toy-H}
\end{equation}
where $x=\sigma-1/2$. Its eigenvalues are
\[
 E_\pm=\pm\sqrt{x^2+|q(t)|^2}.
\]
A zero eigenvalue requires $x=0$ and $q(t)=0$. Thus this Hermitian two-level model separates a horizontal condition from a vertical quantization condition.

The determinant is
\begin{equation}
 \det H(x,t)=-x^2-|q(t)|^2\le0,
 \label{eq:toy-det}
\end{equation}
and vanishes only when both terms vanish. If one independently derived a function $q(t)$ whose zeros were the zeta ordinates and embedded this matrix in a canonical operator family, the conditional spectral theorem would apply.

As written, the toy model is not a zeta construction. Choosing $q(t)=\XiR(t)$ imports the critical-line restriction and does not address off-axis zeros. The exact content of the toy calculation is only the coercive identity in \cref{eq:toy-det}.

\section{Non-Hermitian deformation}

If self-adjointness is broken, let
\[
 H_\epsilon(x,t)=
 \begin{pmatrix}
  x+i\epsilon & q(t)\\
  \overline{q(t)} & -x+i\epsilon
 \end{pmatrix}.
\]
Complex eigenvalues can then occur.

\textbf{INTERPRETACJA.} A candidate physical or dynamical model might describe off-axis displacement as non-Hermitian leakage or failure of positivity. This is heuristic language unless a genuine zeta operator is constructed, and no later theorem depends on this interpretation.

\section{Spectral determinant requirements}

Infinite-dimensional determinants require care. A valid proposal must specify:
\begin{itemize}
  \item the operator domain and proof of self-adjointness;
  \item discreteness or regularization of the relevant spectrum;
  \item the determinant definition and convergence class;
  \item the non-vanishing prefactor relating it to $\XiR$;
  \item multiplicities and possible continuous-spectrum contributions;
  \item arithmetic information explaining the prime side of the explicit formula.
\end{itemize}
Many formal Hamiltonians fail at one of these points.

\section{Random-matrix evidence and its limit}

The local statistics of high zeta zeros resemble eigenvalue statistics of random Hermitian matrices, a connection initiated by Montgomery's pair correlation and developed through extensive numerical and theoretical work \citep{montgomery1973,odlyzko1987,keatingsnaith2000}. Statistical agreement does not identify the operator and does not prove every zero is on the line.

\textbf{INTERPRETACJA.} Within the information-spinor narrative, these statistics motivate treating $t$ as a spectral coordinate. They do not supply evidence for the specific entropy or binary-state map and are not used as a proof step.

\section{Possible integration with the information-spinor theorem}

Suppose, additionally, that a future self-adjoint realization had a swap symmetry and a zero mode decomposed into two sectors, with the symmetry forcing equal sector norms and the eigenvalue equation forcing opposite detector phase. Then the hypotheses of the two-channel cancellation theorem would follow for that zero mode.

\textbf{INTERPRETACJA.} Under those extra assumptions, the information-spinor cancellation picture would be a two-level shadow of the spectral mechanism. This is a conditional reading of a hypothetical construction, not an established property of the zeta function.

\section{Status relative to later chapters}

Later chapters construct exact localized Weil operators under stated hypotheses, an exact positive Riemann kernel, and the quotient entire function $F$. None supplies the spectral determinant required by \cref{thm:spectral-conditional}. The conditional theorem therefore remains valid but uninstantiated, while the later quotient/PF programme proceeds through different proved objects.

\status{The conditional spectral theorem and toy-matrix calculations are exact. Their application to zeta depends on an unconstructed self-adjoint spectral realization. Interpretive readings are explicitly labeled and are not used as proof premises.}
